
\documentclass{article}
\usepackage{colortbl}
\usepackage{makecell}
\usepackage{multirow}
\usepackage{supertabular}

\begin{document}

\newcounter{utterance}

\centering \large Interaction Transcript for game `clean\_up', experiment `1\_medium\_5obj\_et', episode 3 with claude{-}opus{-}4{-}8{-}azure.
\vspace{24pt}

{ \footnotesize  \setcounter{utterance}{1}
\setlength{\tabcolsep}{0pt}
\begin{supertabular}{c@{$\;$}|p{.15\linewidth}@{}p{.15\linewidth}p{.15\linewidth}p{.15\linewidth}p{.15\linewidth}p{.15\linewidth}}
    \# & $\;$A & \multicolumn{4}{c}{Game Master} & $\;\:$B\\
    \hline

    \theutterance \stepcounter{utterance}  
    & & \multicolumn{4}{p{0.6\linewidth}}{
        \cellcolor[rgb]{0.9,0.9,0.9}{
            \makecell[{{p{\linewidth}}}]{
                \texttt{\tiny{[P1$\langle$GM]}}
                \texttt{Ma olen teie mängujuht ja te mängite ühismängu järgmise ruudustikuga mängulauana:} \\
\\ 
\texttt{```} \\
\texttt{ 1234567} \\
\texttt{╔═╤═════╗} \\
\texttt{║H│◌◌◌TI║ 1} \\
\texttt{╟─┤◌◌◌◌◌║ 2} \\
\texttt{║◌└─┐◌◌◌║ 3} \\
\texttt{║◌◌◌├─┐◌║ 4} \\
\texttt{║C◌◌│◌│◌║ 5} \\
\texttt{╟─┬─┴─┼─╢ 6} \\
\texttt{║◌│W◌◌│◌║ 7} \\
\texttt{╚═╧═══╧═╝} \\
\\ 
\texttt{```} \\
\\ 
\texttt{* Ülemisel serval on x{-}koordinaadid, mis suurenevad paremale, ja paremal serval y{-}koordinaadid, mis suurenevad allapoole.} \\
\texttt{* Järgmised objektid on teie ruudustikule juhuslikult paigutatud: 'W', 'I', 'T', 'C', 'H'.} \\
\\ 
\\ 
\texttt{Teine mängija näeb mängulaua varianti, kus objektid on paigutatud erinevatesse juhuslikesse asukohtadesse. Te ei näe teise mängija mängulauda ja tema ei näe teie mängulauda.} \\
\\ 
\\ 
\texttt{**Eesmärk:**} \\
\\ 
\texttt{Mõlemad mängijad peavad oma taustal olevaid objekte liigutama nii, et identsed objektid satuksid samadesse koordinaatidesse. Te peate teise mängijaga suhtlema, et kokku leppida ühises eesmärgi konfiguratsioonis.} \\
\\ 
\\ 
\texttt{**Reeglid:**} \\
\\ 
\texttt{* Igal käigul saate saata täpselt ühe kahest järgmisest käsust:} \\
\texttt{ 1. `ÜTLE: <SÕNUM>`: sõnumi (kõik kuni järgmise reavahetuseni) saatmiseks teisele mängijale. Ma edastan selle teie partnerile.} \\
\texttt{2. `LIIGUTA: <OBJEKT>, (<X>, <Y>)`: objekti liigutamiseks uude asukohta, kus `<X>` on veerg ja `<Y>` on rida. Ma teatan teile, kas teie käik õnnestus või mitte.} \\
\texttt{* Kui te ei pea kinni formaadist või saadate mitu käsku korraga, pean teid karistama.} \\
\texttt{* Kui mõlemad mängijad koguvad rohkem kui 12 karistust, kaotate mõlemad mängu.} \\
\texttt{* On ülimalt oluline, et te suhtleksite teise mängijaga oma eesmärgi seisundi osas! *Ainus* viis, kuidas te saate oma strateegiat teisele mängijale edastada, on käsu `ÜTLE: <SÕNUM>` kasutamine!} \\
\\ 
\\ 
\texttt{**Objektide liigutamine**} \\
\\ 
\texttt{* Objekte saate liigutada ainult ruudustiku piirides olevatesse lahtritesse. Sihtlahter peab olema tühi, st see peab sisaldama ainult sümbolit '◌'.} \\
\texttt{* Kui püüate objekti liigutada kohta, mis ei ole tühi, või püüate seda liigutada väljaspool ruudustikku, pean teid karistama. Saate uue katse.} \\
\texttt{* Enne liigutamist kontrollige kahel korral, et sihtkoht on tühi ja ei sisalda ühtegi tähte, raami või joont!} \\
\\ 
\\ 
\texttt{**Mängu lõpp**} \\
\\ 
\texttt{Kui arvate, et olete saavutanud kõigi objektide joondamise eesmärgi, võite paluda teisel mängijal mäng lõpetada, saates `ÜTLE: valmis?`. Kui teine mängija palub teil mängu lõpetada ja te vastate `ÜTLE: valmis!`, siis mäng lõpeb.} \\
\\ 
\texttt{Mõlemad mängijad võidavad, kui mäng lõpeb 20 vooru jooksul, kus üks voor on määratletud kahe mängija poolt ühe kehtiva käsu saatmisena.} \\
\\ 
\texttt{**Punktiarvestus:**} \\
\\ 
\texttt{Mida lähemal on identsed objektid mõlemal mängulaual, seda rohkem punkte saate. Karistused vähendavad teie punkte. Kas suudate rekordit ületada?} \\
\\ 
\\ 
\texttt{Palun saada teisele mängijale sõnum, et mängu alustada!} \\
            }
        }
    }
    & \\ \\

    \theutterance \stepcounter{utterance}  
    & \multicolumn{4}{p{0.6\linewidth}}{
        \cellcolor[rgb]{0.9,0.9,0.9}{
            \makecell[{{p{\linewidth}}}]{
                \texttt{\tiny{[P1$\rangle$GM]}}
                \texttt{ÜTLE: Tere! Teeme nii, et paneme kõik objektid ritta 2. Ettepanek: W=(3,2), I=(4,2), T=(5,2), C=(6,2), H=(7,2). Kas need lahtrid on sinu laual tühjad?} \\
            }
        }
    }
    & & \\ \\

    \theutterance \stepcounter{utterance}  
    & & & \multicolumn{2}{p{0.3\linewidth}}{
        \cellcolor[rgb]{0.9,0.9,0.9}{
            \makecell[{{p{\linewidth}}}]{
                \texttt{\tiny{[GM$|$GM]}}
                \texttt{ÜTLE: Tere! Teeme nii, et paneme kõik objektid ritta 2. Ettepanek: W=(3,2), I=(4,2), T=(5,2), C=(6,2), H=(7,2). Kas need lahtrid on sinu laual tühjad?} \\
            }
        }
    }
    & & \\ \\

    \theutterance \stepcounter{utterance}  
    & & \multicolumn{4}{p{0.6\linewidth}}{
        \cellcolor[rgb]{0.9,0.9,0.9}{
            \makecell[{{p{\linewidth}}}]{
                \texttt{\tiny{[GM$\rangle$P2]}}
                \texttt{Ma olen teie mängujuht ja te mängite ühismängu järgmise ruudustikuga mängulauana:} \\
\\ 
\texttt{```} \\
\texttt{ 1234567} \\
\texttt{╔═╤═════╗} \\
\texttt{║◌│◌◌◌◌◌║ 1} \\
\texttt{╟─┤◌◌◌H◌║ 2} \\
\texttt{║◌└─┐◌◌◌║ 3} \\
\texttt{║T◌I├─┐◌║ 4} \\
\texttt{║◌◌◌│C│◌║ 5} \\
\texttt{╟─┬─┴─┼─╢ 6} \\
\texttt{║◌│◌◌◌│W║ 7} \\
\texttt{╚═╧═══╧═╝} \\
\\ 
\texttt{```} \\
\\ 
\texttt{* Ülemisel serval on x{-}koordinaadid, mis suurenevad paremale, ja paremal serval y{-}koordinaadid, mis suurenevad allapoole.} \\
\texttt{* Järgmised objektid on teie ruudustikule juhuslikult paigutatud: 'W', 'I', 'T', 'C', 'H'.} \\
\\ 
\\ 
\texttt{Teine mängija näeb mängulaua varianti, kus objektid on paigutatud erinevatesse juhuslikesse asukohtadesse. Te ei näe teise mängija mängulauda ja tema ei näe teie mängulauda.} \\
\\ 
\\ 
\texttt{**Eesmärk:**} \\
\\ 
\texttt{Mõlemad mängijad peavad oma taustal olevaid objekte liigutama nii, et identsed objektid satuksid samadesse koordinaatidesse. Te peate teise mängijaga suhtlema, et kokku leppida ühises eesmärgi konfiguratsioonis.} \\
\\ 
\\ 
\texttt{**Reeglid:**} \\
\\ 
\texttt{* Igal käigul saate saata täpselt ühe kahest järgmisest käsust:} \\
\texttt{ 1. `ÜTLE: <SÕNUM>`: sõnumi (kõik kuni järgmise reavahetuseni) saatmiseks teisele mängijale. Ma edastan selle teie partnerile.} \\
\texttt{2. `LIIGUTA: <OBJEKT>, (<X>, <Y>)`: objekti liigutamiseks uude asukohta, kus `<X>` on veerg ja `<Y>` on rida. Ma teatan teile, kas teie käik õnnestus või mitte.} \\
\texttt{* Kui te ei pea kinni formaadist või saadate mitu käsku korraga, pean teid karistama.} \\
\texttt{* Kui mõlemad mängijad koguvad rohkem kui 12 karistust, kaotate mõlemad mängu.} \\
\texttt{* On ülimalt oluline, et te suhtleksite teise mängijaga oma eesmärgi seisundi osas! *Ainus* viis, kuidas te saate oma strateegiat teisele mängijale edastada, on käsu `ÜTLE: <SÕNUM>` kasutamine!} \\
\\ 
\\ 
\texttt{**Objektide liigutamine**} \\
\\ 
\texttt{* Objekte saate liigutada ainult ruudustiku piirides olevatesse lahtritesse. Sihtlahter peab olema tühi, st see peab sisaldama ainult sümbolit '◌'.} \\
\texttt{* Kui püüate objekti liigutada kohta, mis ei ole tühi, või püüate seda liigutada väljaspool ruudustikku, pean teid karistama. Saate uue katse.} \\
\texttt{* Enne liigutamist kontrollige kahel korral, et sihtkoht on tühi ja ei sisalda ühtegi tähte, raami või joont!} \\
\\ 
\\ 
\texttt{**Mängu lõpp**} \\
\\ 
\texttt{Kui arvate, et olete saavutanud kõigi objektide joondamise eesmärgi, võite paluda teisel mängijal mäng lõpetada, saates `ÜTLE: valmis?`. Kui teine mängija palub teil mängu lõpetada ja te vastate `ÜTLE: valmis!`, siis mäng lõpeb.} \\
\\ 
\texttt{Mõlemad mängijad võidavad, kui mäng lõpeb 20 vooru jooksul, kus üks voor on määratletud kahe mängija poolt ühe kehtiva käsu saatmisena.} \\
\\ 
\texttt{**Punktiarvestus:**} \\
\\ 
\texttt{Mida lähemal on identsed objektid mõlemal mängulaual, seda rohkem punkte saate. Karistused vähendavad teie punkte. Kas suudate rekordit ületada?} \\
\\ 
\texttt{Teine mängija alustas mängu selle sõnumi saatmisega:} \\
\\ 
\texttt{"Tere! Teeme nii, et paneme kõik objektid ritta 2. Ettepanek: W=(3,2), I=(4,2), T=(5,2), C=(6,2), H=(7,2). Kas need lahtrid on sinu laual tühjad?"} \\
\\ 
\texttt{Mis on sinu esimene käsk?} \\
            }
        }
    }
    & \\ \\

    \theutterance \stepcounter{utterance}  
    & & & \multicolumn{4}{p{0.6\linewidth}}{
        \cellcolor[rgb]{0.9,0.9,0.9}{
            \makecell[{{p{\linewidth}}}]{
                \texttt{\tiny{[GM$\langle$P2]}}
                \texttt{Vaatan üle oma laua sihtlahtrid real 2:} \\
\\ 
\texttt{{-} (3,2) = ◌ (tühi) ✓} \\
\texttt{{-} (4,2) = ◌ (tühi) ✓} \\
\texttt{{-} (5,2) = ◌ (tühi) ✓} \\
\texttt{{-} (6,2) = **H** (hõivatud minu H{-}tähega) ⚠️} \\
\texttt{{-} (7,2) = ◌ (tühi) ✓} \\
\\ 
\texttt{Ainuke probleem on, et (6,2) peal on juba minu H. Aga see laheneb kergesti: liigutan H sinna, kuhu see kuulub (7,2), ja siis panen C kohta (6,2).} \\
\\ 
\texttt{**ÜTLE: Sobib! Nõustun: W=(3,2), I=(4,2), T=(5,2), C=(6,2), H=(7,2). Minu laual on (6,2) peal hetkel H, aga liigutan selle esmalt õigesse kohta (7,2) ja teised järele. Alustan!**} \\
            }
        }
    }
     \\ \\

    \theutterance \stepcounter{utterance}  
    & & & \multicolumn{2}{p{0.3\linewidth}}{
        \cellcolor[rgb]{0.9,0.9,0.9}{
            \makecell[{{p{\linewidth}}}]{
                \texttt{\tiny{[GM$|$GM]}}
                \texttt{Vaatan üle oma laua sihtlahtrid real 2:} \\
\\ 
\texttt{{-} (3,2) = ◌ (tühi) ✓} \\
\texttt{{-} (4,2) = ◌ (tühi) ✓} \\
\texttt{{-} (5,2) = ◌ (tühi) ✓} \\
\texttt{{-} (6,2) = **H** (hõivatud minu H{-}tähega) ⚠️} \\
\texttt{{-} (7,2) = ◌ (tühi) ✓} \\
\\ 
\texttt{Ainuke probleem on, et (6,2) peal on juba minu H. Aga see laheneb kergesti: liigutan H sinna, kuhu see kuulub (7,2), ja siis panen C kohta (6,2).} \\
\\ 
\texttt{**ÜTLE: Sobib! Nõustun: W=(3,2), I=(4,2), T=(5,2), C=(6,2), H=(7,2). Minu laual on (6,2) peal hetkel H, aga liigutan selle esmalt õigesse kohta (7,2) ja teised järele. Alustan!**} \\
            }
        }
    }
    & & \\ \\

    \theutterance \stepcounter{utterance}  
    & & & \multicolumn{2}{p{0.3\linewidth}}{
        \cellcolor[rgb]{0.9,0.9,0.9}{
            \makecell[{{p{\linewidth}}}]{
                \texttt{\tiny{[GM$|$GM]}}
                \texttt{Invalid format: head is not empty: 'Vaatan üle oma laua sihtlahtrid real 2:} \\
\\ 
\texttt{{-} (3,2) = ◌ (tühi) ✓} \\
\texttt{{-} (4,2) = ◌ (tühi) ✓} \\
\texttt{{-} (5,2) = ◌ (tühi) ✓} \\
\texttt{{-} (6,2) = **H** (hõivatud minu H{-}tähega) ⚠️} \\
\texttt{{-} (7,2) = ◌ (tühi) ✓} \\
\\ 
\texttt{Ainuke probleem on, et (6,2) peal on juba minu H. Aga see laheneb kergesti: liigutan H sinna, kuhu see kuulub (7,2), ja siis panen C kohta (6,2).} \\
\\ 
\texttt{**'} \\
            }
        }
    }
    & & \\ \\

    \theutterance \stepcounter{utterance}  
    & & \multicolumn{4}{p{0.6\linewidth}}{
        \cellcolor[rgb]{0.9,0.9,0.9}{
            \makecell[{{p{\linewidth}}}]{
                \texttt{\tiny{[GM$\rangle$P2]}}
                \texttt{Karistus: Teie sõnum ei tohi sisaldada midagi enne käsku!} \\
\\ 
\texttt{Veendu, et sinu vastus sisaldab ainult kas `ÜTLE: <SÕNUM>` või `LIIGUTA: <OBJEKT>, (<X>, <Y>)` ja mitte midagi muud!} \\
\\ 
\texttt{Teil on kollektiivselt kogunenud 1 karistust 12 karistusest.} \\
\texttt{Palun proovi uuesti!} \\
            }
        }
    }
    & \\ \\

    \theutterance \stepcounter{utterance}  
    & & & \multicolumn{4}{p{0.6\linewidth}}{
        \cellcolor[rgb]{0.9,0.9,0.9}{
            \makecell[{{p{\linewidth}}}]{
                \texttt{\tiny{[GM$\langle$P2]}}
                \texttt{ÜTLE: Sobib! Nõustun: W=(3,2), I=(4,2), T=(5,2), C=(6,2), H=(7,2). Alustan liigutamist.} \\
            }
        }
    }
     \\ \\

    \theutterance \stepcounter{utterance}  
    & & & \multicolumn{2}{p{0.3\linewidth}}{
        \cellcolor[rgb]{0.9,0.9,0.9}{
            \makecell[{{p{\linewidth}}}]{
                \texttt{\tiny{[GM$|$GM]}}
                \texttt{ÜTLE: Sobib! Nõustun: W=(3,2), I=(4,2), T=(5,2), C=(6,2), H=(7,2). Alustan liigutamist.} \\
            }
        }
    }
    & & \\ \\

    \theutterance \stepcounter{utterance}  
    & & \multicolumn{4}{p{0.6\linewidth}}{
        \cellcolor[rgb]{0.9,0.9,0.9}{
            \makecell[{{p{\linewidth}}}]{
                \texttt{\tiny{[P1$\langle$GM]}}
                \texttt{Sinu sõnum on teisele mängijale edastatud.} \\
\\ 
\\ 
\texttt{Sa mängid praegu vooru 1 maksimaalsest 20 voorust.} \\
\texttt{Teil on kollektiivselt kogunenud 1 karistust 12 karistusest.} \\
\texttt{Teine mängija saatis selle sõnumi:} \\
\\ 
\\ 
\texttt{"Sobib! Nõustun: W=(3,2), I=(4,2), T=(5,2), C=(6,2), H=(7,2). Alustan liigutamist."} \\
\\ 
\\ 
\texttt{Mis on sinu järgmine käsk?} \\
            }
        }
    }
    & \\ \\

    \theutterance \stepcounter{utterance}  
    & \multicolumn{4}{p{0.6\linewidth}}{
        \cellcolor[rgb]{0.9,0.9,0.9}{
            \makecell[{{p{\linewidth}}}]{
                \texttt{\tiny{[P1$\rangle$GM]}}
                \texttt{LIIGUTA: W, (3, 2)} \\
            }
        }
    }
    & & \\ \\

    \theutterance \stepcounter{utterance}  
    & & & \multicolumn{2}{p{0.3\linewidth}}{
        \cellcolor[rgb]{0.9,0.9,0.9}{
            \makecell[{{p{\linewidth}}}]{
                \texttt{\tiny{[GM$|$GM]}}
                \texttt{LIIGUTA: W, (3, 2)} \\
            }
        }
    }
    & & \\ \\

    \theutterance \stepcounter{utterance}  
    & & & \multicolumn{2}{p{0.3\linewidth}}{
        \cellcolor[rgb]{0.9,0.9,0.9}{
            \makecell[{{p{\linewidth}}}]{
                \texttt{\tiny{[GM$|$GM]}}
                \texttt{W moved to (3, 2)} \\
            }
        }
    }
    & & \\ \\

    \theutterance \stepcounter{utterance}  
    & & \multicolumn{4}{p{0.6\linewidth}}{
        \cellcolor[rgb]{0.9,0.9,0.9}{
            \makecell[{{p{\linewidth}}}]{
                \texttt{\tiny{[GM$\rangle$P2]}}
                \texttt{Sinu sõnum on teisele mängijale edastatud.} \\
\\ 
\\ 
\texttt{Sa mängid praegu vooru 2 maksimaalsest 20 voorust.} \\
\texttt{Teil on kollektiivselt kogunenud 1 karistust 12 karistusest.} \\
\texttt{Teine mängija liigutas objekti oma ruudustikul.} \\
\\ 
\\ 
\texttt{Mis on sinu järgmine käsk?} \\
            }
        }
    }
    & \\ \\

    \theutterance \stepcounter{utterance}  
    & & & \multicolumn{4}{p{0.6\linewidth}}{
        \cellcolor[rgb]{0.9,0.9,0.9}{
            \makecell[{{p{\linewidth}}}]{
                \texttt{\tiny{[GM$\langle$P2]}}
                \texttt{LIIGUTA: H, (7, 2)} \\
            }
        }
    }
     \\ \\

    \theutterance \stepcounter{utterance}  
    & & & \multicolumn{2}{p{0.3\linewidth}}{
        \cellcolor[rgb]{0.9,0.9,0.9}{
            \makecell[{{p{\linewidth}}}]{
                \texttt{\tiny{[GM$|$GM]}}
                \texttt{LIIGUTA: H, (7, 2)} \\
            }
        }
    }
    & & \\ \\

    \theutterance \stepcounter{utterance}  
    & & & \multicolumn{2}{p{0.3\linewidth}}{
        \cellcolor[rgb]{0.9,0.9,0.9}{
            \makecell[{{p{\linewidth}}}]{
                \texttt{\tiny{[GM$|$GM]}}
                \texttt{H moved to (7, 2)} \\
            }
        }
    }
    & & \\ \\

    \theutterance \stepcounter{utterance}  
    & & \multicolumn{4}{p{0.6\linewidth}}{
        \cellcolor[rgb]{0.9,0.9,0.9}{
            \makecell[{{p{\linewidth}}}]{
                \texttt{\tiny{[P1$\langle$GM]}}
                \texttt{Edukalt teisaldatud 'W' asukohta (3,2). Teie uuendatud ruudustik näeb välja selline:} \\
\texttt{```} \\
\texttt{ 1234567} \\
\texttt{╔═╤═════╗} \\
\texttt{║H│◌◌◌TI║ 1} \\
\texttt{╟─┤W◌◌◌◌║ 2} \\
\texttt{║◌└─┐◌◌◌║ 3} \\
\texttt{║◌◌◌├─┐◌║ 4} \\
\texttt{║C◌◌│◌│◌║ 5} \\
\texttt{╟─┬─┴─┼─╢ 6} \\
\texttt{║◌│◌◌◌│◌║ 7} \\
\texttt{╚═╧═══╧═╝} \\
\texttt{```} \\
\texttt{Sa mängid praegu vooru 2 maksimaalsest 20 voorust.} \\
\texttt{Teil on kollektiivselt kogunenud 1 karistust 12 karistusest.} \\
\texttt{Teine mängija liigutas objekti oma ruudustikul.} \\
\\ 
\\ 
\texttt{Mis on sinu järgmine käsk?} \\
            }
        }
    }
    & \\ \\

    \theutterance \stepcounter{utterance}  
    & \multicolumn{4}{p{0.6\linewidth}}{
        \cellcolor[rgb]{0.9,0.9,0.9}{
            \makecell[{{p{\linewidth}}}]{
                \texttt{\tiny{[P1$\rangle$GM]}}
                \texttt{LIIGUTA: I, (4, 2)} \\
            }
        }
    }
    & & \\ \\

    \theutterance \stepcounter{utterance}  
    & & & \multicolumn{2}{p{0.3\linewidth}}{
        \cellcolor[rgb]{0.9,0.9,0.9}{
            \makecell[{{p{\linewidth}}}]{
                \texttt{\tiny{[GM$|$GM]}}
                \texttt{LIIGUTA: I, (4, 2)} \\
            }
        }
    }
    & & \\ \\

    \theutterance \stepcounter{utterance}  
    & & & \multicolumn{2}{p{0.3\linewidth}}{
        \cellcolor[rgb]{0.9,0.9,0.9}{
            \makecell[{{p{\linewidth}}}]{
                \texttt{\tiny{[GM$|$GM]}}
                \texttt{I moved to (4, 2)} \\
            }
        }
    }
    & & \\ \\

    \theutterance \stepcounter{utterance}  
    & & \multicolumn{4}{p{0.6\linewidth}}{
        \cellcolor[rgb]{0.9,0.9,0.9}{
            \makecell[{{p{\linewidth}}}]{
                \texttt{\tiny{[GM$\rangle$P2]}}
                \texttt{Edukalt teisaldatud 'H' asukohta (7,2). Teie uuendatud ruudustik näeb välja selline:} \\
\texttt{```} \\
\texttt{ 1234567} \\
\texttt{╔═╤═════╗} \\
\texttt{║◌│◌◌◌◌◌║ 1} \\
\texttt{╟─┤◌◌◌◌H║ 2} \\
\texttt{║◌└─┐◌◌◌║ 3} \\
\texttt{║T◌I├─┐◌║ 4} \\
\texttt{║◌◌◌│C│◌║ 5} \\
\texttt{╟─┬─┴─┼─╢ 6} \\
\texttt{║◌│◌◌◌│W║ 7} \\
\texttt{╚═╧═══╧═╝} \\
\texttt{```} \\
\texttt{Sa mängid praegu vooru 3 maksimaalsest 20 voorust.} \\
\texttt{Teil on kollektiivselt kogunenud 1 karistust 12 karistusest.} \\
\texttt{Teine mängija liigutas objekti oma ruudustikul.} \\
\\ 
\\ 
\texttt{Mis on sinu järgmine käsk?} \\
            }
        }
    }
    & \\ \\

    \theutterance \stepcounter{utterance}  
    & & & \multicolumn{4}{p{0.6\linewidth}}{
        \cellcolor[rgb]{0.9,0.9,0.9}{
            \makecell[{{p{\linewidth}}}]{
                \texttt{\tiny{[GM$\langle$P2]}}
                \texttt{LIIGUTA: C, (6, 2)} \\
            }
        }
    }
     \\ \\

    \theutterance \stepcounter{utterance}  
    & & & \multicolumn{2}{p{0.3\linewidth}}{
        \cellcolor[rgb]{0.9,0.9,0.9}{
            \makecell[{{p{\linewidth}}}]{
                \texttt{\tiny{[GM$|$GM]}}
                \texttt{LIIGUTA: C, (6, 2)} \\
            }
        }
    }
    & & \\ \\

    \theutterance \stepcounter{utterance}  
    & & & \multicolumn{2}{p{0.3\linewidth}}{
        \cellcolor[rgb]{0.9,0.9,0.9}{
            \makecell[{{p{\linewidth}}}]{
                \texttt{\tiny{[GM$|$GM]}}
                \texttt{C moved to (6, 2)} \\
            }
        }
    }
    & & \\ \\

    \theutterance \stepcounter{utterance}  
    & & \multicolumn{4}{p{0.6\linewidth}}{
        \cellcolor[rgb]{0.9,0.9,0.9}{
            \makecell[{{p{\linewidth}}}]{
                \texttt{\tiny{[P1$\langle$GM]}}
                \texttt{Edukalt teisaldatud 'I' asukohta (4,2). Teie uuendatud ruudustik näeb välja selline:} \\
\texttt{```} \\
\texttt{ 1234567} \\
\texttt{╔═╤═════╗} \\
\texttt{║H│◌◌◌T◌║ 1} \\
\texttt{╟─┤WI◌◌◌║ 2} \\
\texttt{║◌└─┐◌◌◌║ 3} \\
\texttt{║◌◌◌├─┐◌║ 4} \\
\texttt{║C◌◌│◌│◌║ 5} \\
\texttt{╟─┬─┴─┼─╢ 6} \\
\texttt{║◌│◌◌◌│◌║ 7} \\
\texttt{╚═╧═══╧═╝} \\
\texttt{```} \\
\texttt{Sa mängid praegu vooru 3 maksimaalsest 20 voorust.} \\
\texttt{Teil on kollektiivselt kogunenud 1 karistust 12 karistusest.} \\
\texttt{Teine mängija liigutas objekti oma ruudustikul.} \\
\\ 
\\ 
\texttt{Mis on sinu järgmine käsk?} \\
            }
        }
    }
    & \\ \\

    \theutterance \stepcounter{utterance}  
    & \multicolumn{4}{p{0.6\linewidth}}{
        \cellcolor[rgb]{0.9,0.9,0.9}{
            \makecell[{{p{\linewidth}}}]{
                \texttt{\tiny{[P1$\rangle$GM]}}
                \texttt{LIIGUTA: T, (5, 2)} \\
            }
        }
    }
    & & \\ \\

    \theutterance \stepcounter{utterance}  
    & & & \multicolumn{2}{p{0.3\linewidth}}{
        \cellcolor[rgb]{0.9,0.9,0.9}{
            \makecell[{{p{\linewidth}}}]{
                \texttt{\tiny{[GM$|$GM]}}
                \texttt{LIIGUTA: T, (5, 2)} \\
            }
        }
    }
    & & \\ \\

    \theutterance \stepcounter{utterance}  
    & & & \multicolumn{2}{p{0.3\linewidth}}{
        \cellcolor[rgb]{0.9,0.9,0.9}{
            \makecell[{{p{\linewidth}}}]{
                \texttt{\tiny{[GM$|$GM]}}
                \texttt{T moved to (5, 2)} \\
            }
        }
    }
    & & \\ \\

    \theutterance \stepcounter{utterance}  
    & & \multicolumn{4}{p{0.6\linewidth}}{
        \cellcolor[rgb]{0.9,0.9,0.9}{
            \makecell[{{p{\linewidth}}}]{
                \texttt{\tiny{[GM$\rangle$P2]}}
                \texttt{Edukalt teisaldatud 'C' asukohta (6,2). Teie uuendatud ruudustik näeb välja selline:} \\
\texttt{```} \\
\texttt{ 1234567} \\
\texttt{╔═╤═════╗} \\
\texttt{║◌│◌◌◌◌◌║ 1} \\
\texttt{╟─┤◌◌◌CH║ 2} \\
\texttt{║◌└─┐◌◌◌║ 3} \\
\texttt{║T◌I├─┐◌║ 4} \\
\texttt{║◌◌◌│◌│◌║ 5} \\
\texttt{╟─┬─┴─┼─╢ 6} \\
\texttt{║◌│◌◌◌│W║ 7} \\
\texttt{╚═╧═══╧═╝} \\
\texttt{```} \\
\texttt{Sa mängid praegu vooru 4 maksimaalsest 20 voorust.} \\
\texttt{Teil on kollektiivselt kogunenud 1 karistust 12 karistusest.} \\
\texttt{Teine mängija liigutas objekti oma ruudustikul.} \\
\\ 
\\ 
\texttt{Mis on sinu järgmine käsk?} \\
            }
        }
    }
    & \\ \\

    \theutterance \stepcounter{utterance}  
    & & & \multicolumn{4}{p{0.6\linewidth}}{
        \cellcolor[rgb]{0.9,0.9,0.9}{
            \makecell[{{p{\linewidth}}}]{
                \texttt{\tiny{[GM$\langle$P2]}}
                \texttt{LIIGUTA: W, (3, 2)} \\
            }
        }
    }
     \\ \\

    \theutterance \stepcounter{utterance}  
    & & & \multicolumn{2}{p{0.3\linewidth}}{
        \cellcolor[rgb]{0.9,0.9,0.9}{
            \makecell[{{p{\linewidth}}}]{
                \texttt{\tiny{[GM$|$GM]}}
                \texttt{LIIGUTA: W, (3, 2)} \\
            }
        }
    }
    & & \\ \\

    \theutterance \stepcounter{utterance}  
    & & & \multicolumn{2}{p{0.3\linewidth}}{
        \cellcolor[rgb]{0.9,0.9,0.9}{
            \makecell[{{p{\linewidth}}}]{
                \texttt{\tiny{[GM$|$GM]}}
                \texttt{W moved to (3, 2)} \\
            }
        }
    }
    & & \\ \\

    \theutterance \stepcounter{utterance}  
    & & \multicolumn{4}{p{0.6\linewidth}}{
        \cellcolor[rgb]{0.9,0.9,0.9}{
            \makecell[{{p{\linewidth}}}]{
                \texttt{\tiny{[P1$\langle$GM]}}
                \texttt{Edukalt teisaldatud 'T' asukohta (5,2). Teie uuendatud ruudustik näeb välja selline:} \\
\texttt{```} \\
\texttt{ 1234567} \\
\texttt{╔═╤═════╗} \\
\texttt{║H│◌◌◌◌◌║ 1} \\
\texttt{╟─┤WIT◌◌║ 2} \\
\texttt{║◌└─┐◌◌◌║ 3} \\
\texttt{║◌◌◌├─┐◌║ 4} \\
\texttt{║C◌◌│◌│◌║ 5} \\
\texttt{╟─┬─┴─┼─╢ 6} \\
\texttt{║◌│◌◌◌│◌║ 7} \\
\texttt{╚═╧═══╧═╝} \\
\texttt{```} \\
\texttt{Sa mängid praegu vooru 4 maksimaalsest 20 voorust.} \\
\texttt{Teil on kollektiivselt kogunenud 1 karistust 12 karistusest.} \\
\texttt{Teine mängija liigutas objekti oma ruudustikul.} \\
\\ 
\\ 
\texttt{Mis on sinu järgmine käsk?} \\
            }
        }
    }
    & \\ \\

    \theutterance \stepcounter{utterance}  
    & \multicolumn{4}{p{0.6\linewidth}}{
        \cellcolor[rgb]{0.9,0.9,0.9}{
            \makecell[{{p{\linewidth}}}]{
                \texttt{\tiny{[P1$\rangle$GM]}}
                \texttt{LIIGUTA: C, (6, 2)} \\
            }
        }
    }
    & & \\ \\

    \theutterance \stepcounter{utterance}  
    & & & \multicolumn{2}{p{0.3\linewidth}}{
        \cellcolor[rgb]{0.9,0.9,0.9}{
            \makecell[{{p{\linewidth}}}]{
                \texttt{\tiny{[GM$|$GM]}}
                \texttt{LIIGUTA: C, (6, 2)} \\
            }
        }
    }
    & & \\ \\

    \theutterance \stepcounter{utterance}  
    & & & \multicolumn{2}{p{0.3\linewidth}}{
        \cellcolor[rgb]{0.9,0.9,0.9}{
            \makecell[{{p{\linewidth}}}]{
                \texttt{\tiny{[GM$|$GM]}}
                \texttt{C moved to (6, 2)} \\
            }
        }
    }
    & & \\ \\

    \theutterance \stepcounter{utterance}  
    & & \multicolumn{4}{p{0.6\linewidth}}{
        \cellcolor[rgb]{0.9,0.9,0.9}{
            \makecell[{{p{\linewidth}}}]{
                \texttt{\tiny{[GM$\rangle$P2]}}
                \texttt{Edukalt teisaldatud 'W' asukohta (3,2). Teie uuendatud ruudustik näeb välja selline:} \\
\texttt{```} \\
\texttt{ 1234567} \\
\texttt{╔═╤═════╗} \\
\texttt{║◌│◌◌◌◌◌║ 1} \\
\texttt{╟─┤W◌◌CH║ 2} \\
\texttt{║◌└─┐◌◌◌║ 3} \\
\texttt{║T◌I├─┐◌║ 4} \\
\texttt{║◌◌◌│◌│◌║ 5} \\
\texttt{╟─┬─┴─┼─╢ 6} \\
\texttt{║◌│◌◌◌│◌║ 7} \\
\texttt{╚═╧═══╧═╝} \\
\texttt{```} \\
\texttt{Sa mängid praegu vooru 5 maksimaalsest 20 voorust.} \\
\texttt{Teil on kollektiivselt kogunenud 1 karistust 12 karistusest.} \\
\texttt{Teine mängija liigutas objekti oma ruudustikul.} \\
\\ 
\\ 
\texttt{Mis on sinu järgmine käsk?} \\
            }
        }
    }
    & \\ \\

    \theutterance \stepcounter{utterance}  
    & & & \multicolumn{4}{p{0.6\linewidth}}{
        \cellcolor[rgb]{0.9,0.9,0.9}{
            \makecell[{{p{\linewidth}}}]{
                \texttt{\tiny{[GM$\langle$P2]}}
                \texttt{LIIGUTA: I, (4, 2)} \\
            }
        }
    }
     \\ \\

    \theutterance \stepcounter{utterance}  
    & & & \multicolumn{2}{p{0.3\linewidth}}{
        \cellcolor[rgb]{0.9,0.9,0.9}{
            \makecell[{{p{\linewidth}}}]{
                \texttt{\tiny{[GM$|$GM]}}
                \texttt{LIIGUTA: I, (4, 2)} \\
            }
        }
    }
    & & \\ \\

    \theutterance \stepcounter{utterance}  
    & & & \multicolumn{2}{p{0.3\linewidth}}{
        \cellcolor[rgb]{0.9,0.9,0.9}{
            \makecell[{{p{\linewidth}}}]{
                \texttt{\tiny{[GM$|$GM]}}
                \texttt{I moved to (4, 2)} \\
            }
        }
    }
    & & \\ \\

    \theutterance \stepcounter{utterance}  
    & & \multicolumn{4}{p{0.6\linewidth}}{
        \cellcolor[rgb]{0.9,0.9,0.9}{
            \makecell[{{p{\linewidth}}}]{
                \texttt{\tiny{[P1$\langle$GM]}}
                \texttt{Edukalt teisaldatud 'C' asukohta (6,2). Teie uuendatud ruudustik näeb välja selline:} \\
\texttt{```} \\
\texttt{ 1234567} \\
\texttt{╔═╤═════╗} \\
\texttt{║H│◌◌◌◌◌║ 1} \\
\texttt{╟─┤WITC◌║ 2} \\
\texttt{║◌└─┐◌◌◌║ 3} \\
\texttt{║◌◌◌├─┐◌║ 4} \\
\texttt{║◌◌◌│◌│◌║ 5} \\
\texttt{╟─┬─┴─┼─╢ 6} \\
\texttt{║◌│◌◌◌│◌║ 7} \\
\texttt{╚═╧═══╧═╝} \\
\texttt{```} \\
\texttt{Sa mängid praegu vooru 5 maksimaalsest 20 voorust.} \\
\texttt{Teil on kollektiivselt kogunenud 1 karistust 12 karistusest.} \\
\texttt{Teine mängija liigutas objekti oma ruudustikul.} \\
\\ 
\\ 
\texttt{Mis on sinu järgmine käsk?} \\
            }
        }
    }
    & \\ \\

    \theutterance \stepcounter{utterance}  
    & \multicolumn{4}{p{0.6\linewidth}}{
        \cellcolor[rgb]{0.9,0.9,0.9}{
            \makecell[{{p{\linewidth}}}]{
                \texttt{\tiny{[P1$\rangle$GM]}}
                \texttt{LIIGUTA: H, (7, 2)} \\
            }
        }
    }
    & & \\ \\

    \theutterance \stepcounter{utterance}  
    & & & \multicolumn{2}{p{0.3\linewidth}}{
        \cellcolor[rgb]{0.9,0.9,0.9}{
            \makecell[{{p{\linewidth}}}]{
                \texttt{\tiny{[GM$|$GM]}}
                \texttt{LIIGUTA: H, (7, 2)} \\
            }
        }
    }
    & & \\ \\

    \theutterance \stepcounter{utterance}  
    & & & \multicolumn{2}{p{0.3\linewidth}}{
        \cellcolor[rgb]{0.9,0.9,0.9}{
            \makecell[{{p{\linewidth}}}]{
                \texttt{\tiny{[GM$|$GM]}}
                \texttt{H moved to (7, 2)} \\
            }
        }
    }
    & & \\ \\

    \theutterance \stepcounter{utterance}  
    & & \multicolumn{4}{p{0.6\linewidth}}{
        \cellcolor[rgb]{0.9,0.9,0.9}{
            \makecell[{{p{\linewidth}}}]{
                \texttt{\tiny{[GM$\rangle$P2]}}
                \texttt{Edukalt teisaldatud 'I' asukohta (4,2). Teie uuendatud ruudustik näeb välja selline:} \\
\texttt{```} \\
\texttt{ 1234567} \\
\texttt{╔═╤═════╗} \\
\texttt{║◌│◌◌◌◌◌║ 1} \\
\texttt{╟─┤WI◌CH║ 2} \\
\texttt{║◌└─┐◌◌◌║ 3} \\
\texttt{║T◌◌├─┐◌║ 4} \\
\texttt{║◌◌◌│◌│◌║ 5} \\
\texttt{╟─┬─┴─┼─╢ 6} \\
\texttt{║◌│◌◌◌│◌║ 7} \\
\texttt{╚═╧═══╧═╝} \\
\texttt{```} \\
\texttt{Sa mängid praegu vooru 6 maksimaalsest 20 voorust.} \\
\texttt{Teil on kollektiivselt kogunenud 1 karistust 12 karistusest.} \\
\texttt{Teine mängija liigutas objekti oma ruudustikul.} \\
\\ 
\\ 
\texttt{Mis on sinu järgmine käsk?} \\
            }
        }
    }
    & \\ \\

    \theutterance \stepcounter{utterance}  
    & & & \multicolumn{4}{p{0.6\linewidth}}{
        \cellcolor[rgb]{0.9,0.9,0.9}{
            \makecell[{{p{\linewidth}}}]{
                \texttt{\tiny{[GM$\langle$P2]}}
                \texttt{LIIGUTA: T, (5, 2)} \\
            }
        }
    }
     \\ \\

    \theutterance \stepcounter{utterance}  
    & & & \multicolumn{2}{p{0.3\linewidth}}{
        \cellcolor[rgb]{0.9,0.9,0.9}{
            \makecell[{{p{\linewidth}}}]{
                \texttt{\tiny{[GM$|$GM]}}
                \texttt{LIIGUTA: T, (5, 2)} \\
            }
        }
    }
    & & \\ \\

    \theutterance \stepcounter{utterance}  
    & & & \multicolumn{2}{p{0.3\linewidth}}{
        \cellcolor[rgb]{0.9,0.9,0.9}{
            \makecell[{{p{\linewidth}}}]{
                \texttt{\tiny{[GM$|$GM]}}
                \texttt{T moved to (5, 2)} \\
            }
        }
    }
    & & \\ \\

    \theutterance \stepcounter{utterance}  
    & & \multicolumn{4}{p{0.6\linewidth}}{
        \cellcolor[rgb]{0.9,0.9,0.9}{
            \makecell[{{p{\linewidth}}}]{
                \texttt{\tiny{[P1$\langle$GM]}}
                \texttt{Edukalt teisaldatud 'H' asukohta (7,2). Teie uuendatud ruudustik näeb välja selline:} \\
\texttt{```} \\
\texttt{ 1234567} \\
\texttt{╔═╤═════╗} \\
\texttt{║◌│◌◌◌◌◌║ 1} \\
\texttt{╟─┤WITCH║ 2} \\
\texttt{║◌└─┐◌◌◌║ 3} \\
\texttt{║◌◌◌├─┐◌║ 4} \\
\texttt{║◌◌◌│◌│◌║ 5} \\
\texttt{╟─┬─┴─┼─╢ 6} \\
\texttt{║◌│◌◌◌│◌║ 7} \\
\texttt{╚═╧═══╧═╝} \\
\texttt{```} \\
\texttt{Sa mängid praegu vooru 6 maksimaalsest 20 voorust.} \\
\texttt{Teil on kollektiivselt kogunenud 1 karistust 12 karistusest.} \\
\texttt{Teine mängija liigutas objekti oma ruudustikul.} \\
\\ 
\\ 
\texttt{Mis on sinu järgmine käsk?} \\
            }
        }
    }
    & \\ \\

    \theutterance \stepcounter{utterance}  
    & \multicolumn{4}{p{0.6\linewidth}}{
        \cellcolor[rgb]{0.9,0.9,0.9}{
            \makecell[{{p{\linewidth}}}]{
                \texttt{\tiny{[P1$\rangle$GM]}}
                \texttt{ÜTLE: Kõik minu objektid on paigas: W=(3,2), I=(4,2), T=(5,2), C=(6,2), H=(7,2). Valmis?} \\
            }
        }
    }
    & & \\ \\

    \theutterance \stepcounter{utterance}  
    & & & \multicolumn{2}{p{0.3\linewidth}}{
        \cellcolor[rgb]{0.9,0.9,0.9}{
            \makecell[{{p{\linewidth}}}]{
                \texttt{\tiny{[GM$|$GM]}}
                \texttt{ÜTLE: Kõik minu objektid on paigas: W=(3,2), I=(4,2), T=(5,2), C=(6,2), H=(7,2). Valmis?} \\
            }
        }
    }
    & & \\ \\

    \theutterance \stepcounter{utterance}  
    & & \multicolumn{4}{p{0.6\linewidth}}{
        \cellcolor[rgb]{0.9,0.9,0.9}{
            \makecell[{{p{\linewidth}}}]{
                \texttt{\tiny{[GM$\rangle$P2]}}
                \texttt{Edukalt teisaldatud 'T' asukohta (5,2). Teie uuendatud ruudustik näeb välja selline:} \\
\texttt{```} \\
\texttt{ 1234567} \\
\texttt{╔═╤═════╗} \\
\texttt{║◌│◌◌◌◌◌║ 1} \\
\texttt{╟─┤WITCH║ 2} \\
\texttt{║◌└─┐◌◌◌║ 3} \\
\texttt{║◌◌◌├─┐◌║ 4} \\
\texttt{║◌◌◌│◌│◌║ 5} \\
\texttt{╟─┬─┴─┼─╢ 6} \\
\texttt{║◌│◌◌◌│◌║ 7} \\
\texttt{╚═╧═══╧═╝} \\
\texttt{```} \\
\texttt{Sa mängid praegu vooru 7 maksimaalsest 20 voorust.} \\
\texttt{Teil on kollektiivselt kogunenud 1 karistust 12 karistusest.} \\
\texttt{Teine mängija saatis selle sõnumi:} \\
\\ 
\\ 
\texttt{"Kõik minu objektid on paigas: W=(3,2), I=(4,2), T=(5,2), C=(6,2), H=(7,2). Valmis?"} \\
\\ 
\\ 
\texttt{Mis on sinu järgmine käsk?} \\
            }
        }
    }
    & \\ \\

    \theutterance \stepcounter{utterance}  
    & & & \multicolumn{4}{p{0.6\linewidth}}{
        \cellcolor[rgb]{0.9,0.9,0.9}{
            \makecell[{{p{\linewidth}}}]{
                \texttt{\tiny{[GM$\langle$P2]}}
                \texttt{ÜTLE: valmis!} \\
            }
        }
    }
     \\ \\

    \theutterance \stepcounter{utterance}  
    & & & \multicolumn{2}{p{0.3\linewidth}}{
        \cellcolor[rgb]{0.9,0.9,0.9}{
            \makecell[{{p{\linewidth}}}]{
                \texttt{\tiny{[GM$|$GM]}}
                \texttt{ÜTLE: valmis!} \\
            }
        }
    }
    & & \\ \\

    \theutterance \stepcounter{utterance}  
    & & \multicolumn{4}{p{0.6\linewidth}}{
        \cellcolor[rgb]{0.9,0.9,0.9}{
            \makecell[{{p{\linewidth}}}]{
                \texttt{\tiny{[P1$\langle$GM]}}
                \texttt{Sinu sõnum on teisele mängijale edastatud.} \\
\\ 
\\ 
\texttt{Sa mängid praegu vooru 7 maksimaalsest 20 voorust.} \\
\texttt{Teil on kollektiivselt kogunenud 1 karistust 12 karistusest.} \\
\texttt{Teine mängija saatis selle sõnumi:} \\
\\ 
\\ 
\texttt{"valmis!"} \\
\\ 
\\ 
\texttt{Mis on sinu järgmine käsk?} \\
            }
        }
    }
    & \\ \\

    \theutterance \stepcounter{utterance}  
    & \multicolumn{4}{p{0.6\linewidth}}{
        \cellcolor[rgb]{0.9,0.9,0.9}{
            \makecell[{{p{\linewidth}}}]{
                \texttt{\tiny{[P1$\rangle$GM]}}
                \texttt{ÜTLE: valmis!} \\
            }
        }
    }
    & & \\ \\

    \theutterance \stepcounter{utterance}  
    & & & \multicolumn{2}{p{0.3\linewidth}}{
        \cellcolor[rgb]{0.9,0.9,0.9}{
            \makecell[{{p{\linewidth}}}]{
                \texttt{\tiny{[GM$|$GM]}}
                \texttt{ÜTLE: valmis!} \\
            }
        }
    }
    & & \\ \\

    \theutterance \stepcounter{utterance}  
    & & & \multicolumn{2}{p{0.3\linewidth}}{
        \cellcolor[rgb]{0.9,0.9,0.9}{
            \makecell[{{p{\linewidth}}}]{
                \texttt{\tiny{[GM$|$GM]}}
                \texttt{true} \\
            }
        }
    }
    & & \\ \\

    \theutterance \stepcounter{utterance}  
    & & & \multicolumn{2}{p{0.3\linewidth}}{
        \cellcolor[rgb]{0.9,0.9,0.9}{
            \makecell[{{p{\linewidth}}}]{
                \texttt{\tiny{[GM$|$GM]}}
                \texttt{Initial states:} \\
\texttt{```} \\
\texttt{Player 1:} \\
\texttt{ 1234567} \\
\texttt{╔═╤═════╗} \\
\texttt{║H│◌◌◌TI║ 1} \\
\texttt{╟─┤◌◌◌◌◌║ 2} \\
\texttt{║◌└─┐◌◌◌║ 3} \\
\texttt{║◌◌◌├─┐◌║ 4} \\
\texttt{║C◌◌│◌│◌║ 5} \\
\texttt{╟─┬─┴─┼─╢ 6} \\
\texttt{║◌│W◌◌│◌║ 7} \\
\texttt{╚═╧═══╧═╝} \\
\\ 
\\ 
\texttt{Player 2:} \\
\texttt{ 1234567} \\
\texttt{╔═╤═════╗} \\
\texttt{║◌│◌◌◌◌◌║ 1} \\
\texttt{╟─┤◌◌◌H◌║ 2} \\
\texttt{║◌└─┐◌◌◌║ 3} \\
\texttt{║T◌I├─┐◌║ 4} \\
\texttt{║◌◌◌│C│◌║ 5} \\
\texttt{╟─┬─┴─┼─╢ 6} \\
\texttt{║◌│◌◌◌│W║ 7} \\
\texttt{╚═╧═══╧═╝} \\
\\ 
\texttt{```} \\
            }
        }
    }
    & & \\ \\

    \theutterance \stepcounter{utterance}  
    & & & \multicolumn{2}{p{0.3\linewidth}}{
        \cellcolor[rgb]{0.9,0.9,0.9}{
            \makecell[{{p{\linewidth}}}]{
                \texttt{\tiny{[GM$|$GM]}}
                \texttt{End states:} \\
\texttt{```} \\
\texttt{Player 1:} \\
\texttt{ 1234567} \\
\texttt{╔═╤═════╗} \\
\texttt{║◌│◌◌◌◌◌║ 1} \\
\texttt{╟─┤WITCH║ 2} \\
\texttt{║◌└─┐◌◌◌║ 3} \\
\texttt{║◌◌◌├─┐◌║ 4} \\
\texttt{║◌◌◌│◌│◌║ 5} \\
\texttt{╟─┬─┴─┼─╢ 6} \\
\texttt{║◌│◌◌◌│◌║ 7} \\
\texttt{╚═╧═══╧═╝} \\
\\ 
\\ 
\texttt{Player 2:} \\
\texttt{ 1234567} \\
\texttt{╔═╤═════╗} \\
\texttt{║◌│◌◌◌◌◌║ 1} \\
\texttt{╟─┤WITCH║ 2} \\
\texttt{║◌└─┐◌◌◌║ 3} \\
\texttt{║◌◌◌├─┐◌║ 4} \\
\texttt{║◌◌◌│◌│◌║ 5} \\
\texttt{╟─┬─┴─┼─╢ 6} \\
\texttt{║◌│◌◌◌│◌║ 7} \\
\texttt{╚═╧═══╧═╝} \\
\\ 
\texttt{```} \\
            }
        }
    }
    & & \\ \\

    \theutterance \stepcounter{utterance}  
    & & & \multicolumn{2}{p{0.3\linewidth}}{
        \cellcolor[rgb]{0.9,0.9,0.9}{
            \makecell[{{p{\linewidth}}}]{
                \texttt{\tiny{[GM$|$GM]}}
                \texttt{* success: True} \\
\texttt{* lose: False} \\
\texttt{* aborted: False} \\
\texttt{{-}{-}{-}{-}{-}{-}{-}} \\
\texttt{* Shifts: 9.00} \\
\texttt{* Max Shifts: 8.00} \\
\texttt{* Min Shifts: 4.00} \\
\texttt{* End Distance Sum: 0.00} \\
\texttt{* Init Distance Sum: 23.93} \\
\texttt{* Expected Distance Sum: 20.95} \\
\texttt{* Penalties: 1.00} \\
\texttt{* Max Penalties: 12.00} \\
\texttt{* Rounds: 7.00} \\
\texttt{* Max Rounds: 20.00} \\
\texttt{* Object Count: 5.00} \\
            }
        }
    }
    & & \\ \\

\end{supertabular}
}

\end{document}
